\clearpage
\section{Positive mirabolic fibres and their regular boundary}
\label{sec:f1-limit-boundary}

The affine commutant at a finite field has a physical matrix trace.  To move
away from arithmetic parameters, one needs an algebraic reason that the same
coefficient trace stays positive.  Orbital valencies supply that reason for
every real $q>1$.  At $q=1$ the form becomes degenerate: the full algebraic
specialization remains, while its reference GNS representation retains only
the ordinary symmetric-group Hecke algebra.  The operational boundary is
therefore built from one-sided sections positive at every nearby real
parameter, rather than from values sampled only at prime powers.

\subsection{Orbit labels, valencies and the positive Gram form}

\begin{definition}[Mirabolic orbit poset and based star-trace]
For $w\in S_n$, define a partial order on $\{1,\ldots,n\}$ by
\[
 i\prec_wj\quad\Longleftrightarrow\quad
 i<j\ \text{ and }\ w^{-1}(i)<w^{-1}(j).
\]
For an antichain $A$, set
\[
 \downarrow_w A=\{i:i=a\text{ or }i\prec_wa
                         \text{ for some }a\in A\}.
\]
The orbit label $(w,A)$ denotes the affine orbital whose flag relative
position is $w$ and whose vector has maximal-support antichain $A$; its
adjacency basis element is $T_{w,A}$.  The empty antichain is the zero-vector
Hecke basis.

On the polynomial orbit-basis algebra define the conjugate-linear star by
\[
 T_{w,A}^*=T_{w^{-1},w^{-1}(A)}
\]
and the coefficient trace by
\[
 \tau_q(T_{w,A})=\begin{cases}1,&w=1, A=\varnothing,\\0,&\text{otherwise}.
 \end{cases}
\]
The proposed positive domain is real $q>1$; evaluation at one is retained as
a positive-semidefinite boundary form.
\end{definition}
\begin{scope}
Vector negation is part of the orbital adjoint but does not change support
antichains.  Positivity is not stipulated and is proved below.  No positive
trace is claimed for $0<q<1$.
\end{scope}
\provenance{D1245, D1246}{theory/sidequests/f1-limit/mirabolic-trace.md}{theory/checks/f1\_limit\_check.py (L12); B1,B4 in theory/checks/f1\_limit\_mirabolic\_check.py}{theory/verdicts/f1-limit-adjudication.md}

\begin{theorem}[Antichain classification and orbital valency]
{\normalfont\statusProved}\quad
Let $L=\mathbb F_Q^n$, fix a reference complete flag, and let $w\in S_n$.
The residual stabilizer orbits on $L$ are indexed by antichains $A$ in the
preceding poset.  The full affine orbital $(w,A)$ has exact valency
\[
 k_{w,A}(Q)=Q^{\ell(w)+|\downarrow_wA|-|A|}(Q-1)^{|A|}.
\]
\end{theorem}
\begin{scope}
This holds for every finite field, including characteristic two and extension
fields.  It is an orbit-counting theorem, not a positivity statement at a
non-arithmetic parameter.
\end{scope}
\provenance{F1-MIR-VAL}{theory/sidequests/f1-limit/mirabolic-trace.md}{theory/checks/f1\_limit\_check.py (L12, n at most 3 and Q=2,3)}{theory/verdicts/f1-limit-adjudication.md}
\begin{proof}
The intersection of the two Borels for relative position $w$ contains the
diagonal torus and precisely the shears $I+cE_{ij}$ with $i\prec_wj$.
Allowed shears change only coordinates below an existing support coordinate,
so the maximal nonzero support is an invariant antichain.  Conversely,
diagonal scaling normalizes the coordinates on $A$, and shears from each
$a\in A$ create arbitrary coordinates in its strict downset.  The residual
orbit therefore has $(Q-1)^{|A|}Q^{|\downarrow_wA|-|A|}$ elements.  The
Bruhat cell of $w$ has $Q^{\ell(w)}$ flags, giving the formula.
\end{proof}

\begin{theorem}[Faithful positive mirabolic trace for all real parameters above one]
{\normalfont\statusProved}\quad
For every $n\geq1$ and every real $q>1$, the orbit-basis trace has diagonal
Gram form
\[
 \tau_q(T_{w,A}^*T_{v,C})
 =\delta_{(w,A),(v,C)}
   q^{\ell(w)+|\downarrow_wA|-|A|}(q-1)^{|A|}.
\]
It is a faithful normalized trace, and its left regular representation makes
the fibre a finite C*-algebra.
\end{theorem}
\begin{scope}
The all-real conclusion follows from a polynomial identity and the displayed
sign, not from numerical sampling at arithmetic values.  For $0<q<1$, odd
$|A|$ gives a negative Gram entry.
\end{scope}
\provenance{F1-MIR-POS}{theory/sidequests/f1-limit/mirabolic-trace.md}{L12 checks finite orbit counts; the uniform positivity statement is proved analytically}{theory/verdicts/f1-limit-adjudication.md}
\begin{proof}
At every prime power, distinct orbital adjacency matrices have disjoint
support, and their normalized Hilbert--Schmidt squared norms equal their
valencies.  Both the product coefficient and the valency are polynomials in
$q$; equality at infinitely many prime powers makes them identical.  For
$q>1$ every displayed diagonal entry is strictly positive.  Polynomial
extension of matrix-trace cyclicity gives traciality.  Left multiplication is
faithful because $L_x1=x$, and its adjoint is $L_{x^*}$.
\end{proof}

\subsection{Hecke retraction and the trace-supported endpoint}

\begin{definition}[Hecke inclusion and vector expectation]
Let
\[
 i_q:H_n(q)\longrightarrow R_n(q),\qquad T_w\longmapsto T_{w,\varnothing},
\]
and define coefficient deletion by
\[
 E_q(T_{w,A})=
 \begin{cases}T_w,&A=\varnothing,\\0,&A\neq\varnothing.
 \end{cases}
\]
\end{definition}
\begin{scope}
These are fixed-rank maps.  They do not provide a generic cross-rank map
$R_m(q)\otimes R_n(q)\to R_{m+n}(q)$; in particular they leave room for
separately constructed right module actions.
\end{scope}
\provenance{D1247}{theory/sidequests/f1-limit/mirabolic-trace.md}{theory/checks/f1\_limit\_mirabolic\_check.py (B4)}{theory/verdicts/f1-limit-adjudication.md}

\begin{theorem}[Vector coarse graining as a conditional expectation]
{\normalfont\statusProved}\quad
For every real $q>1$, $i_q$ is a trace-preserving unital *-monomorphism and
$E_q$ is its trace-preserving UCP conditional expectation, with
$E_qi_q=\mathrm{id}$.  Hence Hecke densities, effects and retained Kraus
lists lift unchanged to the mirabolic fibre.
\end{theorem}
\begin{scope}
The reverse operation $i_qE_q$ is generally a proper vector-sector coarse
graining.  The statement concerns fixed-rank operational data only.
\end{scope}
\provenance{F1-MIR-EXPECT}{theory/sidequests/f1-limit/mirabolic-trace.md}{theory/checks/f1\_limit\_mirabolic\_check.py (B4)}{theory/verdicts/f1-limit-adjudication.md}
\begin{proof}
The Gram formula makes $E_q$ the trace-orthogonal projection onto the unital
Hecke *-subalgebra.  Traciality gives bimodularity.  Testing a positive input
against elements of the subalgebra proves positivity, and the identical
argument in matrix amplifications proves complete positivity.  The basis
formulas give trace preservation and $E_qi_q=\mathrm{id}$.
\end{proof}

\begin{definition}[Algebraic and trace-supported endpoints]
The algebraic endpoint is the full based specialization $R_n(1)$ with the
specialized trace $\tau_1$.  Its trace-supported endpoint is
\[
 R_n(1)/\mathcal N_{\tau_1},\qquad
 \mathcal N_{\tau_1}=\{x:\tau_1(x^*x)=0\}.
\]
The two endpoint algebras are not identified by definition.
\end{definition}
\begin{scope}
No C*-norm is placed on the degenerate algebraic endpoint.  Only its GNS
quotient carries the faithful reference representation.
\end{scope}
\provenance{D1248}{theory/sidequests/f1-limit/mirabolic-trace.md}{theory/checks/f1\_limit\_mirabolic\_check.py (B4); theory/checks/f1\_limit\_check.py (L10)}{theory/verdicts/f1-limit-adjudication.md}

\begin{theorem}[The mirabolic reference quotient at one]
{\normalfont\statusProved}\quad
At $q=1$, the coefficient trace is positive semidefinite and
\[
 \mathcal N_{\tau_1}=\operatorname{span}
       \{T_{w,A}:A\neq\varnothing\}.
\]
This nullspace is a two-sided *-ideal, and
\[
 R_n(1)/\mathcal N_{\tau_1}\cong H_n(1)=\mathbb C[S_n]
\]
as a traced *-algebra.  The map $E_1$ is exactly the quotient map under the
empty-orbit basis identification.
\end{theorem}
\begin{scope}
The full algebraic specialization still contains the vector ideal.  The
claim concerns the GNS quotient selected by the reference trace.
\end{scope}
\provenance{F1-MIR-END}{theory/sidequests/f1-limit/mirabolic-trace.md}{B4, exact rank-two interpolation and null Gram; L10 rank one}{theory/verdicts/f1-limit-adjudication.md}
\begin{proof}
Setting $q=1$ in the diagonal Gram formula leaves norm one exactly for
$A=\varnothing$ and norm zero otherwise.  The nullspace of a positive
tracial form is a two-sided *-ideal by Cauchy--Schwarz.  The surviving basis
classes multiply by the Hecke rule at one, giving the group algebra and its
coefficient trace.  Since $E_1$ deletes precisely the null basis, it is the
quotient.
\end{proof}

\subsection{The actual one-sided operational category}

\begin{definition}[Regular one-sided labels]
Fix $J_\varepsilon=[1,1+\varepsilon]$.  A regular mirabolic section is
\[
 h(q)=\sum_{w,A}h_{w,A}(q)T_{w,A}(q),
 \qquad h_{w,A}\in C(J_\varepsilon),
\]
with pointwise polynomial multiplication, the orbit star and continuous
coefficient trace.  Hecke factors use their continuous section algebra on
the same interval; separated register words use continuous coefficients in
the product bases and the $C(J_\varepsilon)$-balanced common-parameter
tensor.

A regular density additionally satisfies $h(q)\geq0$ and
$\tau_q(h(q))=1$ for every real $1<q\leq1+\varepsilon$.  A regular effect
satisfies $0\leq e(q)\leq1$ there; a regular POVM consists of finitely many
such positive sections summing to one.  Regular instrument coefficients are
continuous and obey Kraus completeness at every real $q>1$ in the interval.
\end{definition}
\begin{scope}
There is no C*-norm on the boundary section algebra.  Positivity only at
prime powers is insufficient.  Coefficient-pole densities are excluded and
require separate singular boundary data.
\end{scope}
\provenance{D1249}{theory/sidequests/f1-limit/mirabolic-boundary.md}{uniform theorem by proof; L10 and B4 give boundary examples}{theory/verdicts/f1-limit-adjudication.md}

\begin{definition}[Regular separated operational category]
For each real $q>1$, quantum atoms are $[H_n]$ and $[R_n]$, $n\geq1$, with
their faithful coefficient traces; nonempty finite classical wires have
uniform trace.  The empty word has algebra $\mathbb C$, and ordered word
tensor is the actual spatial tensor product.  The interval version uses the
regular common-parameter sections just defined.

Raw morphisms are finite typed planar circuits generated by all regular
preparations, discards and finite POVMs; every retained same-register
instrument
\[
 K:X\to X\underline O,
 \qquad \sum_{o,i}K_{o,i}(q)^*K_{o,i}(q)=1;
\]
the explicit classical product, singleton, bijective relabeling and
classical-wire routing maps; and
\[
 \begin{aligned}
 \operatorname{up}_n&:[H_n]\to[R_n]&&\quad(E_q\text{ in Heisenberg form}),\\
 \operatorname{down}_n&:[R_n]\to[H_n]&&\quad(i_q\text{ in Heisenberg form}).
 \end{aligned}
\]
Instrument lists are retained under separated context tensor and factorwise
inclusion.  With the uniform classical trace, the $o$-block of the output
density is $|O|\sum_iK_{o,i}hK_{o,i}^*$.

Composition is typed gluing and ordered tensor is horizontal juxtaposition.
Equality is the congruence generated by category, monoidal and explicit
classical-wiring relations, literal state/POVM equality, pointwise equality
of coefficient Kraus Grams within each outcome, the routed list identities,
and $\operatorname{down}_n\operatorname{up}_n=\mathrm{id}_{[H_n]}$.
The reverse composite is $i_qE_q$ and is not an identity.  Germs identify
circuits when these relations hold on a common smaller one-sided interval.
There is no cross-rank mirabolic assembly generator.
\end{definition}
\begin{scope}
The classical wiring is the explicit parameter-independent UCP presentation
already proved for the core category; no unresolved wiring hypothesis remains.
Equality is not an isolated-CP quotient and requires no continuous family of
Kraus mixing matrices.  In particular $[R_m][R_n]$ is not identified with
$[R_{m+n}]$.
\end{scope}
\provenance{D1257, D1218}{theory/sidequests/f1-limit/mirabolic-boundary.md; theory/sidequests/f1-limit/classical-wiring.md}{theory/checks/f1\_limit\_wiring\_check.py (W1--W5); B4}{theory/verdicts/f1-limit-adjudication.md}

\begin{definition}[Regular quotient boundary functor]
Let $\Pi_n:R_n(1)\to H_n(1)$ be the GNS quotient $E_1$.  For a register word
$X$, let $Q_X$ tensor $\Pi_n$ on every mirabolic factor and the identity on
Hecke and classical factors; write $\overline X$ for the word obtained by
replacing every mirabolic atom by its Hecke atom.  The endpoint category has
these Hecke/classical words, the same wiring, all normalized preparations,
POVMs and retained same-register instruments, the same label relations, and
its UCP realization.

The boundary functor evaluates continuous coefficients at one, applies $Q_X$
to every density, effect and Kraus entry, and sends every up/down generator
to the identity.  Evaluation at $q>1$ requires a chosen representative whose
interval contains $q$.
\end{definition}
\begin{scope}
Arbitrary UCP families with merely regular superoperator coefficients are not
source morphisms.  The regular density and retained-Kraus restrictions are
part of the functor's domain.
\end{scope}
\provenance{D1258}{theory/sidequests/f1-limit/mirabolic-boundary.md}{structured descent proof; B4 and W1--W5 test primitive data}{theory/verdicts/f1-limit-adjudication.md}

\begin{theorem}[Regular operational mirabolic boundary]
{\normalfont\statusProved}\quad
The regular one-sided circuits and their germs form actual ordered monoidal
categories with UCP realizations.  Coefficient evaluation followed by the
factorwise GNS quotient defines a monoidal endpoint functor to the separated
Hecke operational category.  It preserves all specified generators and
typed relations.  Finite protocol Born weights converge, and conditioned
weights converge whenever the limiting conditioning probability is positive.
Every individual finite endpoint Hecke circuit has a local regular lift.
\end{theorem}
\begin{scope}
All source densities, effects and instruments are positive or normalized for
every nearby real $q>1$.  Arithmetic-only positivity and arbitrary
coefficient-regular UCP superoperators are excluded.  No generic
$R_m(q)\otimes R_n(q)$ assembly is asserted.
\end{scope}
\provenance{F1-MIR-REG}{theory/sidequests/f1-limit/mirabolic-boundary.md}{B4 and W1--W5 support primitive formulas; functoriality and positivity by structured proof}{theory/verdicts/f1-limit-adjudication.md}
\begin{proof}
Polynomial structure constants and the parameter-independent star make
regular sections a *-algebra.  For a regular positive density $h$ and any
endpoint Hecke element $x$, lift $x$ through the Hecke inclusion.  Positivity
of $\tau_q(x(q)^*h(q)x(q))$ for every real $q>1$ passes to one by coefficient
continuity.  Faithfulness of the endpoint quotient trace then forces
$Q_X(h(1))\geq0$; normalization also passes to the limit.  The same argument
handles effects and POVMs, while Kraus completeness passes algebraically and
multiplicativity of $Q_X$ makes every branch map descend.

The explicit wiring commutes with factorwise quotients.  Pointwise Gram
equality is a congruence fibre by fibre and descends under the linear quotient
map.  Hence the circuit quotient and endpoint functor are well defined.  A
finite branch density has continuous coefficients by induction over its
circuit, so trace pairings and finite event sums converge.  Positive-limit
conditioning is ordinary quotient continuity.  Local endpoint lifts use
square-root normalization in the continuous Hecke subalgebra followed by
$i_q$, on one common neighborhood for the finitely many labels.
\end{proof}

The real-neighborhood hypothesis is sharp.  In $H_2(q)$, if $e_+,e_-$ are
the two spectral projections, then
\[
 h(q)=(q-	frac32)e_++\frac{5}{2q}e_-
\]
is normalized and positive at every prime power $Q\geq2$, but has a negative
$e_+$ coefficient for $1<q<3/2$.  It is not a regular density.  Conversely,
the rank-one singular state $q(q-1)^{-1}(1-e)$ is positive and normalized for
every $q>1$, but has a coefficient pole and does not annihilate the endpoint
null ideal.  It is a legitimate punctured-fibre state requiring additional
boundary data, not an object of the regular category.
